% GENERATED FILE. Edit canonical lessons, then run npm run generate:books.
% canonical-source-hash: fnv1a64:a1dc7b470f614430
% canonical-lessons: LA-C553-vestibulum, LA-C553-pessulus, LA-C553-limen, LA-C553-cardo, LA-C553-penus

\chapter{Doors and stores}
\label{ch:la-doors-and-stores}
\hlchaptermodality{553}

\begin{chapteropening}
\textbf{By the end of this chapter:} \emph{I can describe the entrance of a house and its stores.}

Complete the last lesson of chapter 553: I can describe the entrance of a house and its stores.
\end{chapteropening}

\section[vestibulum, vestibulī]{vestibulum, vestibulī — a forecourt, entrance court}

\label{lesson:LA-C553-vestibulum}



\begin{quote}
Before the new one: say the Latin for to sweep, then the Latin for a broom.
\end{quote}

\subsection*{You'll want to know: vestibulum, vestibulī}
\textbf{vestibulum, vestibulī} — \textquotedblleft{}a forecourt, entrance court\textquotedblright{}.

\begin{quote}
Hospitēs in vestibulō exspectant.
\end{quote}

\begin{quote}
\emph{The guests wait in the forecourt.}
\end{quote}

Neuter, second declension. The space between the street and the front door. English \emph{vestibule}.

\begin{cousinweb}
\textbf{vestibulum} is of uncertain origin; Lewis \& Short connects it with a root meaning \textquotedblleft{}to dwell\textquotedblright{}, but it is not settled. English \emph{vestibule} comes from it.
\end{cousinweb}

\begin{grammarlens}[title={What you've built}]
One more word for telling, explaining and describing at B1.
\end{grammarlens}

\subsection*{Your turn}
\begin{itemize}
\raggedright
  \item \emph{Say it:} \emph{vestibulum, vestibulī}
  \item \emph{Say it:} \emph{vestibulum, vestibulī}, once more
  \item \emph{Say it:} a sentence of your own with \emph{vestibulum, vestibulī}
  \item \emph{Recall:} say the Latin for to sweep, then the Latin for a broom, then say \emph{vestibulum, vestibulī} again
\end{itemize}

\begin{tcolorbox}[breakable,colback=teal!4,colframe=teal!35!black,title={Before you move on}]
What does vestibulum, vestibulī mean? (\textquotedblleft{}A forecourt, entrance court\textquotedblright{}.) Say it once more.
\end{tcolorbox}
\section[pessulus, pessulī]{pessulus, pessulī — a door bolt}

\label{lesson:LA-C553-pessulus}



\begin{quote}
Before the new one: say the Latin for a broom, then the Latin for a forecourt.
\end{quote}

\subsection*{You'll want to know: pessulus, pessulī}
\textbf{pessulus, pessulī} — \textquotedblleft{}a door bolt\textquotedblright{}.

\begin{quote}
Iānitor noctū pessulum iānuae obdit.
\end{quote}

\begin{quote}
\emph{At night the doorkeeper shoots the bolt of the door.}
\end{quote}

Masculine, second declension. \emph{Pessulum obdere} — to shoot the bolt.

\begin{cousinweb}
\textbf{pessulus} was borrowed from Greek \emph{passalos} \textquotedblleft{}a peg\textquotedblright{} (Lewis \& Short; EtymDB).
\end{cousinweb}

\begin{grammarlens}[title={What you've built}]
One more word for telling, explaining and describing at B1.
\end{grammarlens}

\subsection*{Your turn}
\begin{itemize}
\raggedright
  \item \emph{Say it:} \emph{pessulus, pessulī}
  \item \emph{Say it:} \emph{pessulus, pessulī}, once more
  \item \emph{Say it:} a sentence of your own with \emph{pessulus, pessulī}
  \item \emph{Recall:} say the Latin for a broom, then the Latin for a forecourt, then say \emph{pessulus, pessulī} again
\end{itemize}

\begin{tcolorbox}[breakable,colback=teal!4,colframe=teal!35!black,title={Before you move on}]
What does pessulus, pessulī mean? (\textquotedblleft{}A door bolt\textquotedblright{}.) Say it once more.
\end{tcolorbox}
\section[līmen, līminis]{līmen, līminis — a threshold, doorstep}

\label{lesson:LA-C553-limen}



\begin{quote}
Before the new one: say the Latin for a forecourt, then the Latin for a door bolt.
\end{quote}

\subsection*{You'll want to know: līmen, līminis}
\textbf{līmen, līminis} — \textquotedblleft{}a threshold, doorstep\textquotedblright{}.

\begin{quote}
Nova nūpta līmen nōn tangit.
\end{quote}

\begin{quote}
\emph{The new bride does not touch the threshold.}
\end{quote}

Neuter, third declension. A bride was carried over the threshold, so as not to stumble on it. English \emph{subliminal}.

\begin{cousinweb}
\textbf{līmen} is a native Latin noun; Lewis \& Short connects it with \emph{līmus} \textquotedblleft{}sidelong\textquotedblright{}, a crossbeam, but the origin is not settled. English \emph{subliminal} comes from it.
\end{cousinweb}

\begin{grammarlens}[title={What you've built}]
One more word for telling, explaining and describing at B1.
\end{grammarlens}

\subsection*{Your turn}
\begin{itemize}
\raggedright
  \item \emph{Say it:} \emph{līmen, līminis}
  \item \emph{Say it:} \emph{līmen, līminis}, once more
  \item \emph{Say it:} a sentence of your own with \emph{līmen, līminis}
  \item \emph{Recall:} say the Latin for a forecourt, then the Latin for a door bolt, then say \emph{līmen, līminis} again
\end{itemize}

\begin{tcolorbox}[breakable,colback=teal!4,colframe=teal!35!black,title={Before you move on}]
What does līmen, līminis mean? (\textquotedblleft{}A threshold, doorstep\textquotedblright{}.) Say it once more.
\end{tcolorbox}
\section[cardō, cardinis]{cardō, cardinis — a hinge}

\label{lesson:LA-C553-cardo}



\begin{quote}
Before the new one: say the Latin for a door bolt, then the Latin for a threshold.
\end{quote}

\subsection*{You'll want to know: cardō, cardinis}
\textbf{cardō, cardinis} — \textquotedblleft{}a hinge\textquotedblright{}.

\begin{quote}
Iānua vetus in cardine sonat.
\end{quote}

\begin{quote}
\emph{The old door creaks on its hinge.}
\end{quote}

Masculine, third declension. Also a turning point, and so English \emph{cardinal} — that on which things turn.

\begin{cousinweb}
\textbf{cardō} is a native Latin noun of uncertain origin; Lewis \& Short offers only guesses. English \emph{cardinal}, \textquotedblleft{}chief, pivotal\textquotedblright{}, comes from it.
\end{cousinweb}

\begin{grammarlens}[title={What you've built}]
One more word for telling, explaining and describing at B1.
\end{grammarlens}

\subsection*{Your turn}
\begin{itemize}
\raggedright
  \item \emph{Say it:} \emph{cardō, cardinis}
  \item \emph{Say it:} \emph{cardō, cardinis}, once more
  \item \emph{Say it:} a sentence of your own with \emph{cardō, cardinis}
  \item \emph{Recall:} say the Latin for a door bolt, then the Latin for a threshold, then say \emph{cardō, cardinis} again
\end{itemize}

\begin{tcolorbox}[breakable,colback=teal!4,colframe=teal!35!black,title={Before you move on}]
What does cardō, cardinis mean? (\textquotedblleft{}A hinge\textquotedblright{}.) Say it once more.
\end{tcolorbox}
\section[penus, penūs]{penus, penūs — provisions, the store of food}

\label{lesson:LA-C553-penus}



\begin{quote}
Before the new one: say the Latin for a threshold, then the Latin for a hinge.
\end{quote}

\subsection*{You'll want to know: penus, penūs}
\textbf{penus, penūs} — \textquotedblleft{}provisions, the store of food\textquotedblright{}.

\begin{quote}
Vīlica penum in cellā servat.
\end{quote}

\begin{quote}
\emph{The housekeeper keeps the provisions in the storeroom.}
\end{quote}

Masculine or feminine, fourth declension; also \emph{penum}, neuter. The stored food of a household.

\begin{cousinweb}
\textbf{penus} is inherited from a Proto-Indo-European word for \textquotedblleft{}food\textquotedblright{} (EtymDB).
\end{cousinweb}

\begin{grammarlens}[title={What you've built}]
One more word for telling, explaining and describing at B1.
\end{grammarlens}

\subsection*{Your turn}
\begin{itemize}
\raggedright
  \item \emph{Say it:} \emph{penus, penūs}
  \item \emph{Say it:} \emph{penus, penūs}, once more
  \item \emph{Say it:} a sentence of your own with \emph{penus, penūs}
  \item \emph{Recall:} say the Latin for a threshold, then the Latin for a hinge, then say \emph{penus, penūs} again
\end{itemize}

\begin{tcolorbox}[breakable,colback=teal!4,colframe=teal!35!black,title={Before you move on}]
What does penus, penūs mean? (\textquotedblleft{}Provisions, the store of food\textquotedblright{}.) Say it once more.
\end{tcolorbox}
